\slajdid{08-s044}
\begin{frame}[label=08-s044]
\gusttekst
\begin{columns}[c,onlytextwidth]\column{.40\textwidth}\centering\input{slike/tikz/s042_rc_struje.tex}
\column{.58\textwidth}
\begin{align*}
u_i(t)-u_R(t)-u_C(t)&=0\\
i_R(t)&=i_C(t)\\[5pt]
u_R(t)+u_C(t)&=u_i(t)\\
Ri_R(t)+u_C(t)&=u_i(t)\\
Ri_C(t)+u_C(t)&=u_i(t)\\
RC\frac{du_C(t)}{dt}+u_C(t)&=u_i(t)
\end{align*}
\[u_i(t)=(U-U_0)h(t-t_0)+U_0\]
gde je h(t) Hejvisajdova funkcija. Kako je
\[u_o(t)=u_C(t)\]
\[RC\frac{du_o(t)}{dt}+u_o(t)=u_i(t)\]
\end{columns}
\end{frame}
